%! Observational Studies — Unit 1, Lesson 1.5 — Guided Notes (Answer Key)
\documentclass[10pt]{article}
\usepackage{saar-article}
\usepackage{saar-key}   % pulls in -boxes; do NOT also load -boxes
\newcommand{\TallMath}[1]{$\displaystyle #1\rule[-1.4em]{0pt}{3.2em}$}

\begin{document}

\pageheader{Unit 1, Lesson 1.5}{Guided Notes --- Answer Key}
% NAMESTRIP: no \namedateperiod here — mirrors the blank. Teacher-only prose goes
% in the lesson plan as \begin{teachernote}[Guided Notes], never in a key.

\vspace{0.15in}

\begin{objectivebox}
\small
By the end of this lesson, I will be able to\dots
\begin{enumerate}[leftmargin=1.5em, itemsep=2pt, topsep=2pt]
  \item explain what makes a study \textbf{observational} --- nobody is assigned
        to a group;
  \item use the \textbf{statistical cycle} to plan an observational study that
        answers a question;
  \item compare two groups that already exist and describe the
        \textbf{association} between two variables;
  \item explain why an association \textbf{does not} show causation, and name a
        \textbf{lurking variable} that could explain it instead;
  \item decide which conclusions an observational study \emph{does} and
        \emph{does not} support, and justify the decision.
\end{enumerate}
\end{objectivebox}

\boxguard
\begin{vocabbox}
\small
Fill in each term as we build it together.
\vocabans{Observational study}{A study in which you measure or survey the
members of a sample without trying to affect them. The researcher does not
assign anyone to a group --- the groups already exist. A survey is one type of
observational study.}
\vocabans{Explanatory variable}{The variable we think might explain or account
for the other one --- at the pool, the time of day a member swims. It is the
variable at the tail of the arrow in the claim being tested.}
\vocabans{Response variable}{The variable we measure to see what happens --- at
the pool, whether a member gets at least $7$ hours of sleep. It sits at the head
of the arrow.}
\vocabans{Association}{Two variables are associated when they move together:
knowing one group a member is in changes what you expect for the other
variable. An association is what an observational study can show.}
\vocabans{Causation}{One variable actually produces the change in the other.
Causation is a much stronger claim than association, and watching alone is never
enough evidence for it.}
\vocabans{Lurking variable}{A variable outside the study that affects both of
the variables in it, and can create an association all by itself --- at the
pool, whether a member still works full time.}
\end{vocabbox}

\boxguard[12]
\begin{hookbox}
\small
Harbor Point's director looks at a clean, unbiased survey of $200$ members:
$70\%$ of the morning swimmers get at least $7$ hours of sleep, but only $45\%$
of the evening swimmers do. She writes a flyer: \emph{``Swim in the morning ---
sleep better!''}

\vspace{2pt}
\textbf{The sample was random, nobody was left out, and every count is correct.
Is her flyer's claim safe? Why or why not?}
\ansline{Answers vary. No --- the members picked their own swim time, so the two groups}
\ansline{differ in other ways too. The survey shows the two go together, not that one causes it.}
\end{hookbox}

\boxguard[30]
\begin{notesbox}{1. What Makes a Study Observational}
\small
In an \textbf{observational study} you measure or survey the members of a sample
\emph{without trying to affect them}. You watch what is already happening. The
researcher \textbf{does not assign} anybody to a group --- the groups already
exist, because people put themselves there.

\vspace{3pt}
Which of the director's plans are observational studies?

\begin{center}
\renewcommand{\arraystretch}{1.4}
\begin{tabularx}{0.98\linewidth}{@{} X c p{3.4cm} @{}}
\toprule
\textbf{The plan} & \textbf{Observational?} & \textbf{Who decided the groups?} \\
\midrule
Ask a random sample of members what time they usually swim and how much they
sleep. & \ans{yes} & \ans{the members did} \\
Count how many members walk in before $8$~a.m.\ every day for a week.
  & \ans{yes} & \ans{nobody --- just watching} \\
Assign half the members to swim at $6$~a.m.\ and half at $7$~p.m.\ for a month,
then compare their sleep. & \ans{no} & \ans{the director did} \\
Mail the sleep survey to every member on the membership list.
  & \ans{yes} & \ans{the members did} \\
\bottomrule
\end{tabularx}
\end{center}

\vspace{3pt}
A \textbf{survey} is one type of observational study --- which means
\emph{every} study in Unit 1 so far has been one.

\vspace{2pt}
\textit{The whole difference is one word: \emph{assign}. If the researcher put
people into the groups, it is not an observational study.}
\end{notesbox}

\boxguard[30]
\begin{notesbox}{2. Planning an Observational Study --- the Statistical Cycle}
\small
Lesson 1.1 gave you the four stages. Here is the pool study, one stage at a
time. Fill in what \emph{this} study does at each stage.

\begin{center}
\renewcommand{\arraystretch}{1.45}
\begin{tabularx}{0.98\linewidth}{@{} p{2.6cm} p{4.4cm} X @{}}
\toprule
\textbf{Stage} & \textbf{What the stage asks} & \textbf{Our pool study} \\
\midrule
Formulate & What do we want to know, and about whom?
  & Do members who swim in the morning report more sleep than members who swim
    in the evening? \\
Collect & Who is on the frame, and how is the sample chosen?
  & \ans{a random sample of $200$ from the full $1{,}500$-member list} \\
Represent & How will we organize what comes back?
  & \ans{a two-way table of swim time against sleep, plus a percent for each
    group} \\
Analyze \& communicate & What do the numbers say, \emph{in context}?
  & \ans{compare $70\%$ to $45\%$ and report the $25$-point gap} \\
\bottomrule
\end{tabularx}
\end{center}

\vspace{3pt}
Notice what is \emph{not} anywhere in this plan: nowhere do we tell a single
member when to swim. That one absence is what makes it \ans{an observational study}.

\vspace{2pt}
\textit{The cycle is the same cycle. What changes from study to study is what
you are allowed to say at the end of it.}
\end{notesbox}

\boxguard[30]
\begin{notesbox}{3. Comparing Two Groups That Already Exist}
\small
All $200$ members answered. Here is the whole study in one table:

\begin{center}
\renewcommand{\arraystretch}{1.4}
\begin{tabularx}{0.95\linewidth}{@{} X c c c @{}}
\toprule
 & \textbf{At least 7 hrs} & \textbf{Fewer than 7} & \textbf{Total} \\
\midrule
Morning swimmers & $56$ & $24$ & $80$ \\
Evening swimmers & $54$ & $66$ & $120$ \\
\midrule
Total & $110$ & $90$ & $200$ \\
\bottomrule
\end{tabularx}
\end{center}

\vspace{3pt}
From your warm-up: \ans{$70\%$} of morning swimmers and \ans{$45\%$} of
evening swimmers get at least $7$ hours --- a gap of \ans{$25$} percentage
points, with the \ans{morning} group higher.

\vspace{3pt}
About how many of all $1{,}500$ members get at least $7$ hours of sleep?

\begin{work}
  \dfrac{110}{200} &= 0.55 \\
  0.55 \times 1500 &= 825 \text{ members}
\end{work}

\vspace{2pt}
Here the \textbf{explanatory variable} (the one we think might explain the
other) is \ans{swim time}, and the \textbf{response variable} (the one we measure
to see what happens) is \ans{hours of sleep}.

\vspace{2pt}
\textit{Nobody built these two groups. The members chose their own swim time
long before anyone asked them a question.}
\end{notesbox}

\boxguard[30]
\begin{notesbox}{4. Association Is Not Causation}
\small
The $25$-point gap is real, and it is not bias --- the sample was random and
everyone answered. Two variables that move together like this are
\textbf{associated}. But there are at least three different stories that would
produce exactly this table:

\begin{enumerate}[leftmargin=1.6em, itemsep=4pt, topsep=3pt]
  \item \textbf{The claim:} swimming in the morning $\rightarrow$ more sleep.
  \item \textbf{The arrow runs backwards:} \ans{people who sleep well can get up at $6$ a.m.}
  \item \textbf{Something else explains both.} The same survey recorded whether
        each member is retired:
\end{enumerate}

\begin{center}
\renewcommand{\arraystretch}{1.4}
\begin{tabularx}{0.9\linewidth}{@{} X c c @{}}
\toprule
 & \textbf{Retired} & \textbf{Percent retired} \\
\midrule
Morning swimmers ($80$) & $60$ & \ans{$75\%$} \\
Evening swimmers ($120$) & $18$ & \ans{$15\%$} \\
\bottomrule
\end{tabularx}
\end{center}

\vspace{3pt}
Working a full-time job changes \emph{when you can swim} and \emph{how much you
sleep}. A variable that sits outside the study and affects both of the variables
in it is a \ans{lurking variable}.

\vspace{3pt}
An observational study can show that two variables are \ans{associated}. It cannot
show that one of them \ans{caused} the other, because nobody assigned the
groups.

\vspace{2pt}
\textit{Association means they travel together. Causation means one of them
does the driving. Watching cannot tell you which.}
\end{notesbox}

\boxguard[30]
\begin{notesbox}{5. What This Study Can and Cannot Claim}
\small
Same data, three statements. Decide which ones the study supports.

\begin{center}
\renewcommand{\arraystretch}{1.45}
\begin{tabularx}{0.98\linewidth}{@{} X c p{3.9cm} @{}}
\toprule
\textbf{Statement} & \textbf{Supported?} & \textbf{Why / why not} \\
\midrule
At Harbor Point, morning swimmers report more sleep than evening swimmers.
  & \ans{yes} & \ans{it just reports what was measured} \\
Swimming in the morning makes a person sleep longer.
  & \ans{no} & \ans{causal claim; nobody was assigned a swim time} \\
If the evening swimmers switched to mornings, they would sleep more.
  & \ans{no} & \ans{predicts a change, so it assumes causation} \\
\bottomrule
\end{tabularx}
\end{center}

\vspace{3pt}
\textbf{Two rules for writing up an observational study.} Fill in the missing
word:

\begin{enumerate}[leftmargin=1.6em, itemsep=4pt, topsep=3pt]
  \item Report what you \emph{saw}: name both groups, both percents, and the
        gap --- and say the two variables are \ans{associated}.
  \item Never use a word like \emph{causes}, \emph{makes}, \emph{raises}, or
        \emph{improves}, and never predict what would happen if someone
        \ans{switched} groups.
\end{enumerate}

\vspace{3pt}
To find out whether the swim time \emph{itself} does anything, somebody would
have to \ans{assign members to a swim time at random} --- which is exactly what next class is about.
\end{notesbox}

\boxguard
\begin{practicebox}
\small
Each row reports a real association found by watching. Name the two variables
and one lurking variable that could explain the association instead.

\begin{center}
\renewcommand{\arraystretch}{1.5}
\begin{tabularx}{0.98\linewidth}{@{} X p{2.5cm} p{2.5cm} p{2.7cm} @{}}
\toprule
\textbf{What was observed} & \textbf{Explanatory} & \textbf{Response}
  & \textbf{A lurking variable} \\
\midrule
Students who eat breakfast score higher on math tests.
  & \ans{eating breakfast} & \ans{math score} & \ans{sleep and home routine} \\
The more firefighters sent to a fire, the more damage the fire does.
  & \ans{firefighters sent} & \ans{damage done} & \ans{size of the fire} \\
Card holders who attend story hour read more books per year.
  & \ans{story hour} & \ans{books read} & \ans{families who already read} \\
\bottomrule
\end{tabularx}
\end{center}

\vspace{3pt}
A news story reports: ``Teens who play a school sport have better grades, so
every school should require a sport.'' Explain, in one or two sentences, why
this observational study does not support that recommendation.
\ansline{Nobody assigned teens to play a sport --- they chose, and the ones who chose may}
\ansline{already have had better grades, more time, or more support. It is an association only.}
\end{practicebox}

\end{document}
